\section{Terminal junctions and the coefficients that can vary}
\label{sec:terminal}

We work over $\mathbb C$.  Lemma~\ref{lem:cycle-rank} reduces a
hypothetical infinite klt program to a small tail.  We lift that tail
to crepant terminal models, joined by coefficient decreases and finite
strings of small flips, and show that every coefficient which keeps
varying belongs to a divisor whose downstairs center has codimension
at least three.

Two finiteness statements make this construction useful.  Finiteness of
the exceptional valuations with log discrepancy at most one lets us
stabilize the divisors extracted on the terminal models.  The description
of valuations below discrepancy two then identifies the ordinary
families that must be subtracted to obtain a finite difficulty.  Both
follow from the same log-smooth calculation.

\subsection{A log-smooth calculation}

We use characteristic-zero resolution and principalization in the
form originating with Hironaka \cite{Hironaka1964I,Hironaka1964II}:
the chosen resolution can make the finitely many divisor supports
and ideal transforms simple normal crossings while preserving the
open set where they already have that form.

\begin{lemma}\label{lem:snc-valuations}
Let $\pi:W\to X$ be a log resolution of a klt pair, and include the exceptional locus and the crepant boundary support in a simple normal crossings (SNC) divisor $\sum_jT_j$. Write $a_j=a(T_j;X,B)>0$, assigning $a_j=1$ to any auxiliary component of coefficient zero. If $J$ indexes the components through the generic center of a divisorial valuation $\xi$ on $W$, then
\begin{equation}\label{eq:snc-formula}
 a(\xi;X,B)=\sum_{j\in J}\xi(T_j)a_j+e(\xi).
\end{equation}
Here $e(\xi)$ is a nonnegative integer, the log discrepancy for the reduced SNC boundary near that center. If the center has codimension $c$ in $W$, then
\[
 e(\xi)\geq c-|J|.
\]
If $e(\xi)=0$, then $\xi$ is the monomial valuation associated to a primitive positive integral vector over the generic point of an SNC stratum. In particular there are finitely many such valuations when their orders $\xi(T_j)$ are bounded.
\end{lemma}

\begin{proof}
Near the center the crepant boundary is
$\sum_{j\in J}(1-a_j)T_j$. Subtracting it from the reduced boundary gives \eqref{eq:snc-formula}. The reduced SNC pair is lc and has integral Cartier adjoint, so its log discrepancies are nonnegative integers. At the generic center choose $c-|J|$ additional regular parameters cutting that center inside its stratum. Adjoining their divisors with coefficient one still gives an SNC pair locally. Log canonicity of this larger reduced boundary implies
$e(\xi)\geq\sum_\ell\xi(t_\ell)\geq c-|J|$.

For completeness, the zero-error assertion can be seen by a toroidal modification. The orders along the crossing equations give a positive integral vector in the cone of the stratum. Make a regular fan subdivision containing its primitive ray; the resulting local toric modification, pulled back through the crossing coordinates, is log crepant for the reduced SNC structure. The lifted center lies in the open stratum of the divisor for this ray. If it were a proper subvariety of that stratum, the preceding transverse-parameter argument would give positive reduced-SNC discrepancy. Thus it is the generic point of that divisor, and the normalized valuation is its order. The toric construction is split over the residue field of the original stratum, so a primitive vector determines one divisor there. There are finitely many strata and finitely many bounded integral vectors.
\end{proof}

\begin{lemma}\label{lem:small-discrepancies}
For a fixed klt pair $(X,B)$ there are finitely many exceptional divisorial valuations with $a(\xi;X,B)\leq1$.
\end{lemma}

\begin{proof}
Choose the resolution of Lemma~\ref{lem:snc-valuations} with SNC exceptional divisor. If the lifted center meets that SNC divisor and contributes discrepancy at most one, positivity of its $a_j$ forces $e(\xi)=0$ and bounds every positive order in \eqref{eq:snc-formula}. There are finitely many resulting monomial valuations. If the center avoids the SNC divisor, an exceptional valuation has center of codimension at least two on the smooth isomorphism open, and hence log discrepancy at least two. Divisors already on the resolution form a finite list of possible exceptionals. These cases exhaust the valuations.
\end{proof}

\begin{lemma}[Echoes below log discrepancy two]\label{lem:echoes}
Let $(Y,\Theta)$ be an effective terminal pair on a projective $\mathbb Q$-factorial fourfold. Label finitely many prime divisors $S_h$ containing $\operatorname{Supp}\Theta$, and write
$\Theta=\sum_hb_hS_h$, allowing $b_h=0$.
Then:
\begin{enumerate}[label=\textup{(\roman*)}]
\item Only finitely many exceptional valuations with $a(\xi;Y,\Theta)<2$ have center of codimension at least three.
\item At each fixed surface center there are only finitely many exceptional valuations with log discrepancy less than two.
\item Apart from finitely many surface centers, every valuation with surface center and log discrepancy less than two is centered on the one-label smooth open. In regular parameters $x,t$ at its generic point, with $x=0$ defining that label $S_h$, the contributing valuations are exactly the monomial orders
\begin{equation}\label{eq:echo-values}
 \xi(x)=j,\quad \xi(t)=1,\qquad
 a(\xi;Y,\Theta)=1+j(1-b_h)<2,\qquad j\geq1.
\end{equation}
\end{enumerate}
More strongly, all exceptional valuations below two which are not of the form \eqref{eq:echo-values} on that one-label smooth open belong to a finite list.
\end{lemma}

\begin{proof}
The underlying variety $Y$ is terminal. Indeed $\Theta$ is effective and $\mathbb Q$-Cartier, so its removal only increases discrepancies. In particular $Y$ is smooth in codimension two. Choose a closed set $A$ of codimension at least two containing $\operatorname{Sing}Y$, the nonsmooth loci of the labels, and every intersection of two distinct labels. On $Y\setminus A$ the ambient space is smooth and at most one smooth label is present.

Resolve and principalize the ideal of $A$, simultaneously with the label divisors, by a morphism isomorphic over $Y\setminus A$. The full inverse image of $A$ has SNC divisorial support, all of whose components are exceptional over $Y$. Their log discrepancies for $(Y,\Theta)$ are strictly greater than one. Any valuation centered in $A$ has its lifted center in one of these exceptional components. If its discrepancy is less than two, \eqref{eq:snc-formula} therefore forces $e(\xi)=0$; all orders are bounded because all coefficients $a_j$ are positive. Lemma~\ref{lem:snc-valuations} gives only finitely many such valuations.

On the smooth open away from all labels, an exceptional valuation has discrepancy at least two. At a center of codimension $c$ in a single label of coefficient $b_h$, the same formula gives
\[
 a(\xi;Y,\Theta)=(1-b_h)\xi(S_h)+e(\xi),
 \qquad e(\xi)\geq c-1.
\]
Since $1-b_h>0$, discrepancy less than two forces $c=2$. At its generic point choose regular parameters $x,t$, with $x$ the label equation. Add $t=0$ as an auxiliary component of coefficient zero to the SNC calculation. Then
\[
 a(\xi;Y,\Theta)=(1-b_h)\xi(x)+\xi(t)+e'(\xi).
\]
The strict bound and positivity force $\xi(t)=1$ and $e'(\xi)=0$. The zero-error characterization gives precisely \eqref{eq:echo-values}. Conversely each such primitive vector defines the indicated divisorial order and discrepancy. Only finitely many $j$ satisfy the displayed inequality. Combining this with the finite list over $A$ proves every assertion.
\end{proof}

The local classification in Lemma~\ref{lem:echoes} is the echo phenomenon underlying the difficulties of \cite[Example~1.4 and Lemma~1.5]{AHK}; their discrepancy convention is one less than ours. We have included the argument in order to specify the finite correction terms used later.

\subsection{Crepant terminalizations and the finite bridges}

Lifting a prescribed program to terminal models with decreasing
exceptional coefficients is also used in
\cite[Lemma~2.8]{HanLiuZhuang2025}. We give the construction here,
including the comparison that identifies the prescribed positive end.

Suppose now that
$(X_i,B_i)\dashrightarrow(X_{i+1},B_{i+1})$
is a small klt tail, with step bases $Z_i$.
By Lemmas~\ref{lem:graph} and \ref{lem:small-discrepancies}, the sets
\[
 I_i=\{\xi:\xi\text{ exceptional over }X_i,\ a(\xi;X_i,B_i)\leq1\}
\]
are nested finite sets: exceptional status is unchanged under a small map. After truncation they equal one fixed set $I_0$. Write $a_i(\xi)=a(\xi;X_i,B_i)$.

\begin{proposition}\label{prop:terminal-chain}
There are projective birational morphisms $p_i:Y_i\to X_i$ extracting exactly $I_0$, with $Y_i$ $\mathbb Q$-factorial and
\begin{equation}\label{eq:terminal-crepant}
 K_{Y_i}+\Theta_i=p_i^*(K_{X_i}+B_i),\qquad
 \Theta_i=B_i^{\mathrm{str}}+
       \sum_{s\in I_0}(1-a_i(s))E_{i,s},
\end{equation}
such that each $(Y_i,\Theta_i)$ is effective terminal. Consecutive junctions can be joined by:
\begin{enumerate}[label=\textup{(\roman*)}]
\item decreasing the coefficient of each $E_{i,s}$ to $1-a_{i+1}(s)$ on $Y_i$;
\item a finite sequence, possibly empty, of small negative terminal flips over $Z_i$.
\end{enumerate}
Every labeled prime persists through the chain. All coefficients are nonincreasing, and the discrepancy of every valuation is nondecreasing, including across coefficient changes.
\end{proposition}

\begin{proof}
For klt input, \cite[Corollary~1.4.3]{BCHM} extracts precisely any prescribed finite set of exceptional valuations with log discrepancies at most one by a projective birational morphism with $\mathbb Q$-factorial source. The restriction in that corollary on discrepancy-one centers is vacuous for a klt pair. Apply it to $I_0$. The boundary in \eqref{eq:terminal-crepant} is effective with coefficients less than one. Any valuation exceptional over $Y_i$ is exceptional over $X_i$ and is not one of the extracted divisors; it therefore has discrepancy greater than one. This proves terminality.

Put
\[
 \Theta'_i=B_i^{\mathrm{str}}+
       \sum_{s\in I_0}(1-a_{i+1}(s))E_{i,s},
 \qquad A_i=K_{Y_i}+\Theta'_i.
\]
The new coefficients lie in $[0,1)$ and are no larger than their old values. Since $Y_i$ is $\mathbb Q$-factorial, the removed effective boundary is $\mathbb Q$-Cartier. Its order at every valuation is nonnegative, so this move preserves terminality and increases discrepancies weakly.

On a smooth common model $W$ of $Y_i$ and $X_{i+1}$ over $Z_i$, with projections $q$ and $r$, consider
\[
 E=q^*A_i-r^*D_{i+1}.
\]
The divisor $A_i$ is the divisorial trace on $Y_i$ of the pullback of $D_{i+1}$. At primes from $X_i$ this follows from smallness downstairs. At each marked prime $E_{i,s}$ it is exactly the definition of its new coefficient using $a_{i+1}(s)$. Thus $q_*E=0$. Also $-E$ is $q$-nef because $D_{i+1}$ is ample over $Z_i$. Negativity gives $E\geq0$. Every prime on $X_{i+1}$ corresponds to a prime on $Y_i$ with the same coefficient, so $E$ is exceptional over $X_{i+1}$ as well.

Run a chosen finite klt MMP for $A_i$ over $Z_i$. The precise existence input is \cite[Theorem~1.2 and Corollary~1.4.2]{BCHM}: a $\mathbb Q$-factorial klt pair with relatively big boundary admits a terminating program with suitable scaling. Here $Y_i\to Z_i$ is birational, so both its boundary and its adjoint are relatively big; their generic fibre is a point, as in \cite[Definition~3.1.1(7)]{BCHM}. Choose a general effective rational scaling divisor with sufficiently positive class and sufficiently small component coefficients, so the augmented pair remains klt and its adjoint is relatively nef. A fibre-type output over $Z_i$ is impossible by dimensions. Hence a finite choice reaches a relatively nef end $(\widehat Y,\widehat\Theta)$.

For this finite program, the input pullback equals the nef-end pullback plus an effective end-exceptional divisor. We already have the comparison with the ample divisor $D_{i+1}$ on $X_{i+1}$. Lemma~\ref{lem:two-ends} therefore gives a morphism
$\widehat p:\widehat Y\to X_{i+1}$ and the actual crepant identity
\[
 K_{\widehat Y}+\widehat\Theta=\widehat p^*D_{i+1}.
\]
No marked divisor can have been contracted during the finite program: its discrepancy at the input after the coefficient decrease is $a_{i+1}(s)$, and crepancy gives the same value at the end, whereas a contraction losing it would give strict improvement by Lemma~\ref{lem:graph}. No other prime can have been lost either. Such a prime comes from $X_i$, hence from $X_{i+1}$, and a birational morphism to the normal $X_{i+1}$ has a strict transform of that prime. The program extracts no divisors, so a lost prime could not reappear.

Every step was consequently small. Terminality persists by Corollary~\ref{cor:singularities-preserved}, and the end again extracts exactly $I_0$. Set it equal to $Y_{i+1}$. This constructs the junctions inductively and proves all discrepancy and coefficient assertions. Figure~\ref{fig:terminal-bridge} summarizes the construction.
\end{proof}

\begin{figure}[ht]
\centering
\begin{tikzcd}[column sep=large,row sep=large]
 (Y_i,\Theta_i) \arrow[r,"\Theta_i\searrow\Theta_i'"]
 \arrow[d,"p_i"']
 & (Y_i,\Theta_i') \arrow[r,dashed,"\text{finite small flips}"{yshift=3pt}]
 & (Y_{i+1},\Theta_{i+1}) \arrow[d,"p_{i+1}"]\\
 (X_i,B_i) \arrow[r,"f_i"']
 & Z_i
 & (X_{i+1},B_{i+1}) \arrow[l,"f_i^+"]
\end{tikzcd}
\caption{The terminal junctions are crepant over the prescribed downstairs models. Coefficients decrease before the finite relative bridge; the original downstairs choices remain fixed.}
\label{fig:terminal-bridge}
\end{figure}

Let $S_h$ label the strict transforms of the original boundary components and the finitely many marked extractions throughout this chain, retaining labels whose coefficient becomes zero. \mbox{Write the current boundary as}
$\Theta=\sum_h b_hS_h$. The coefficient $b_h$ is constant at the original boundary labels and nonincreasing at every extracted label. Call a label \emph{fixed} if its coefficient is eventually constant, and \emph{varying} otherwise. After truncating once more, all fixed coefficients are constant. A varying coefficient is positive at every remaining junction, since reaching zero would make it fixed.

\subsection{Surface mld ACC and the centers of varying labels}

For a codimension-two subvariety $V\subset X_i$, write
\[
 \operatorname{mld}_{\eta_V}(X_i,B_i)
 =\inf\{a_i(\xi):c_{X_i}(\xi)=V\}.
\]
We use two surface facts. Minimal log discrepancies at points of klt surfaces with coefficients in a fixed finite set satisfy ACC; this follows from \cite[Theorems~3.1 and 3.8]{Alexeev1993}. To spell out the fixed-set reduction, omit zero coefficients. If the smallest permitted positive coefficient is $c$, the local coefficient-sum bound $\sum b_j\leq2$ bounds the number of components through the point by $\lfloor2/c\rfloor$. Only finitely many coefficient multisets occur. A strictly increasing sequence of mld values would therefore have a subsequence with one constant coefficient vector, contradicting Theorem~3.8 of that reference. If the positive coefficient set is empty, the vector is already constant and empty. This argument also allows coefficients repeated by splitting on a surface section. A singular codimension-two point of an effective log pair has mld at most one; see \cite[Proposition~3.2, arXiv version]{Ambro1999}.

\begin{lemma}\label{lem:transverse-mld}
The numbers $\operatorname{mld}_{\eta_V}(X_i,B_i)$, as $i$ and $V$ vary, belong to the set of surface mld values with the same finite set of boundary coefficients. Each minimum is attained.
\end{lemma}

\begin{proof}
Fix a model and a center $V$. Choose a projective log resolution which also principalizes the ideal of $V$. Its divisorial inverse image has finitely many components dominating $V$. The mld in question is the minimum of their log discrepancies. Indeed every valuation centered at $\eta_V$ has its lifted center in at least one such component, and the positive-coefficient formula \eqref{eq:snc-formula} bounds its discrepancy below by that component's discrepancy. Conversely those component valuations have center $V$.

Cut the fourfold by two general sufficiently ample hyperplanes so that the resulting surface meets $V$ at general points and avoids the proper closed images in $V$ of resolution strata which do not dominate it. Use simultaneous Bertini for the finitely many relevant strata. The downstairs surface is normal; the two pulled-back cuts on the smooth resolution are smooth and transverse to the SNC data. They contain no exceptional divisor as a component, because neither general hyperplane contains an exceptional image. Their complete intersection has no component contained in the exceptional locus: each exceptional divisor has image of dimension at most two, and two general pulled-back cuts leave dimension at most one there.

Adjunction cancels the two normal determinant factors on the source and target of the restricted resolution. Thus the crepant coefficients on the surface resolution are the restrictions of the original coefficients. The boundary on the normal surface has the original finite coefficient set: generic hyperplane sections of each boundary prime are reduced, and no new boundary divisor is introduced. All sliced crepant coefficients have positive log discrepancy, so the surface pair is klt. Components of the principalized inverse image dominating $V$ cut to divisors with the same discrepancies over each chosen general intersection point. Strata with proper image were avoided. Applying \eqref{eq:snc-formula} on the slice now gives the same minimum as above. This proves the assertion.
\end{proof}

\begin{proposition}\label{prop:codim-two}
After truncating the small downstairs tail, every surface contained in either bad locus of every remaining step has generic-point mld greater than one and is generically in the smooth locus of its ambient model. At each remaining crepant junction, every varying label has downstairs center of codimension at least three. Consequently the reduced union of varying labels is mapped by $p_i$ to a set of dimension at most one.
\end{proposition}

\begin{proof}
Suppose infinitely many steps have a bad surface with mld at most one on one specified side. A computing divisor is exceptional downstairs and belongs to the stabilized finite set $I_0$. Some fixed member $s$ computes the mld at infinitely many such occurrences. Its discrepancies along this subsequence increase strictly. On the negative side, strict improvement occurs in the step immediately following the occurrence; on the positive side, it occurs immediately before the occurrence. Monotonicity between occurrences then gives strict increase in either case. By Lemma~\ref{lem:transverse-mld}, these are surface mld values with a fixed finite coefficient set, contradicting ACC. Apply this to both sides and discard finitely many initial steps. Ambro's bound gives generic ambient smoothness at the remaining bad surfaces.

If an extracted label with discrepancy at most one has a codimension-two center $V$ downstairs at some remaining junction, $V$ cannot be contained in the following negative bad locus. Otherwise its mld would be at most the discrepancy of that label, contradicting the first assertion. The birational step is therefore an isomorphism near the generic point of $V$, and the center and its discrepancy persist on the next model. The same argument repeats at every later step. This label's coefficient is hence constant from that junction onward. A varying label cannot have such a center. It is exceptional downstairs, so its center has codimension at least three; in dimension four its image has dimension at most one.
\end{proof}
